\documentclass[12pt,a4paper]{article}
\usepackage[utf8]{inputenc}
\usepackage[spanish]{babel}
\usepackage{amsmath}
\usepackage{amssymb}
\usepackage{listings}
\usepackage{enumitem}
\usepackage{geometry}
\geometry{margin=2.5cm}
\usepackage[utf8]{inputenc} % allow utf-8 input

\usepackage{listings}
\usepackage{xcolor}

% Define colors
\definecolor{codegreen}{rgb}{0,0.6,0}
\definecolor{codegray}{rgb}{0.5,0.5,0.5}
\definecolor{codepurple}{rgb}{0.58,0,0.82}
\definecolor{backcolour}{rgb}{0.95,0.95,0.92}

% Configure listings
\lstdefinestyle{mystyle}{
    backgroundcolor=\color{backcolour},   
    commentstyle=\color{codegreen},
    keywordstyle=\color{magenta},
    numberstyle=\tiny\color{codegray},
    stringstyle=\color{codepurple},
    basicstyle=\ttfamily\footnotesize,
    breakatwhitespace=false,         
    breaklines=true,                 
    captionpos=b,                    
    keepspaces=true,                 
    numbers=left,                    
    numbersep=5pt,                  
    showspaces=false,                
    showstringspaces=false,
    showtabs=false,                  
    tabsize=2
}

\lstset{style=mystyle}


\setlength{\parindent}{0pt} % Eliminar sangría de primera línea en todos los párrafos
% add space between paragraphs
\setlength{\parskip}{1em}

% agregar carpeta imagenes
\usepackage{graphicx}
\graphicspath{ {../imagenes/} }


\title{Aprendizaje Reforzado: Guía 2: Procesos de decisión Markovianos}
\author{José Saint Germain}

\begin{document}

\maketitle

\section{Ejercicio 1}

Para un proceso de Markov donde $P_{ss'} = P[S_{t+1} = s' | S_t = s]$ 
son las probabilidades de transición, demostrar que la probabilidad de un
episodio se puede escribir como el producto de las probabilidades de transición
y la condición inicial, es decir:
$$P[S_0=s_0, S_1=s_1, ..., S_t=s_t] = P_{s_{t}s_{t+1}} P_{s_{t}s_{t-1}s_t} \dots P_{s_1 s_2}P_{s_0 s_1}P_{s_{0'}}$$

Es decir, demostrar que la probabilidad conjunta de un episodio (serie de estados)
se puede calcular por el producto de las transiciones multiplicado por la probabildad
del estado inicial.

\subsection{Resolución}

Empezamos demostrándolo con la probabilidad conjunta de tres estados:

Por teorema de Bayes sabemos que:
$$ P(S_0,S_1,S_2) = P(S_2|S_1,S_0)P(S_1,S_0) $$

Sin embargo $P(S_2|S_1,S_0) = P(S_2|S_1)$ por la propiedad de Markov, por lo que:

$$ P(S_0,S_1,S_2) = P(S_2|S_1)P(S_1,S_0) $$

Luego, volvemos a descomponer $P(S_1,S_0)$ con el teorema de Bayes:

$$ P(S_0,S_1,S_2) = P(S_2|S_1)P(S_1|S_0)P(S_0) $$

Realizando una inducción, podemos ver que para $n$ estados se cumple:
$$ P(S_0,S_1,...,S_n) = P(S_n|S_{n-1})P(S_{n-1}|S_{n-2})...P(S_1|S_0)P(S_0) $$

\section{Ejercicio 2}

Para el proceso de Markov de la figura, calcular la probabilidad de
cada uno de los siguientes episodios condicionados al estado inicial C1:

\begin{enumerate}
    \item C1, C2, C3, Pass, Sleep
    \item C1, FB, FB, C1, C2, Sleep
    \item C1, C2, C3, Pub, C2, C3, Pass, Sleep
    \item C1, FB, FB, C1, C2, C3, Pub, C1, FB, FB, FB, C1, C2, C3, Pub, C2, Sleep
\end{enumerate}

\subsection{Resolución}

\subsubsection{Inciso 1}

\begin{align}
    \tag{1} &= P(C1,C2,C3,Pass,Sleep|C1) \\
    \tag{2} &= P(C2|C1)P(C3|C2)P(Pass|C3)P(Sleep|Pass)\\
    \tag{4} &= 0.5 * 0.8 * 0.6 * 1 = 0.24
\end{align}

\subsubsection{Inciso 2}

\begin{align}
\tag{1} &= P(C1,FB,FB,C1,C2,Sleep|C1) \\
\tag{2} &= P(FB|C1)P(FB|FB)P(C1|FB)P(C2|C1)P(Sleep|C2) \\
\tag{4} &= 0.5 * 0.9 * 0.1 * 0.5 * 0.2 = 0.0045 
\end{align}

\subsubsection{Inciso 3}

\begin{align}
\tag{1} &= P(C1,C2,C3,Pub,C2,C3,Pass,Sleep|C1) \\
\tag{2} &= P(C2|C1)P(C3|C2)P(Pub|C3)P(C2|Pub)P(C3|C2)P(Pass|C3)P(Sleep|Pass) \\
\tag{4} &= 0.5 * 0.8 * 0.4 * 0.4 * 0.8 * 0.6 * 1 = 0.03073
\end{align}

\subsection{Inciso 4}

\begin{align}
\tag{1} &= P(C1,FB,FB,C1,C2,C3,\dots ,Sleep|C1) \\
\tag{2} &= P(FB|C1)P(FB|FB)P(C1|FB)P(C2|C1)P(C3|C2),\dots,P(Sleep|C2) \\
\tag{4} &= .5 * .9 * .1 * .5 * .8 * .4 * .2 * .5 * .9 * .9 * .1 * .5 * .8 * .4 * .4 * .2 = 7.46e^{-07}
\end{align}

\section{Ejercicio 3}

Demostrar que si la recompensa es constante $R_t = R \forall t$, 
y el factor de descuento es $\gamma<1$, entonces $$G_t = \frac{R}{1-\gamma}$$

\subsection{Resolución}

Para resolver el ejercicio recordamos la definición de una serie geométrica:
$$ S= \sum_{i=0}^n q^n = \frac{1-q^{n+1}}{1-q} $$
y que si $|q|<1$ entonces:
$$ S= \sum_{i=0}^\infty q^n = \frac{1}{1-q} $$

Entonces, si el retorno es:
$$ G_t = \sum_{k=0}^\infty \gamma^k R_{t+k+1} $$
y $R_t$ es constante, podemos reescribir y utilizar la serie geométrica:
$$ G_t = \sum_{k=0}^\infty \gamma^k R = R \sum_{k=0}^\infty \gamma^k \stackrel{\text{serie geom}}{=} R \frac{1}{1-\gamma} = \frac{R}{1-\gamma} $$

\section{Ejercicio 4}

Para el proceso de Recompensas Markoviano (MRP) de la figura,
calcular el retorno $G_0$ de los episodios del ejercicio 2
con estado inicial C1 y $\gamma=0.5$.

\subsection{Resolución}

\subsubsection{Inciso 1}

\begin{align}
\tag{1} &= R_{C2} + \gamma^1 R_{C3} + \gamma^2 R_{PASS}  + \gamma^3 R_{SLEEP} \\
\tag{2} &= -2      + 0.5 * -2      + 0.5^2 * 10         + 0.5^3 * 0 \\
\tag{3} &= -2      -1       + 2.5            + 0 = -0.5
\end{align}

\subsubsection{Inciso 2}
\begin{align}
\tag{1} &= R_{FB} + \gamma R_{FB} + \gamma^2 R_{C1} + \gamma^3 R_{C2} + \gamma^4 R_{SLEEP} \\
\tag{2} &= -1      + 0.5 * -1      + 0.5^2 * -2      + 0.5^3 * -2      + 0.5^4 * 0 \\
\tag{3} &= -1      -0.5            -0.5              -0.25             + 0 = -2.25
\end{align}

\subsubsection{Inciso 3}
\begin{align}
\tag{1} &= R_{C2} + \gamma R_{C3} + \gamma^2 R_{PUB} + \gamma^3 R_{C2} + \gamma^4 R_{C3} + \gamma^5 R_{PASS} + \gamma^6 R_{SLEEP} \\
\tag{2} &= -2      + 0.5 * -2      + 0.5^2 * 1      + 0.5^3 * -2      + 0.5^4 * -2      + 0.5^5 * 10         + 0.5^6 * 0 \\
\tag{3} &= -2      -1              + 0.25           -0.25             -0.125            + 0.3125          + 0 = -2.8125
\end{align}

\subsubsection{Inciso 4}
\begin{align}
\tag{1} &= R_{FB} + \gamma R_{FB} + \gamma^2 R_{C1} + \dots + \gamma^{15} R_{SLEEP} \\
\tag{2} &= -1      + 0.5 * -1      + 0.5^2 * -2      + \dots + 0.5^{15} * 0 \\
\tag{3} &= -1      -0.5            -0.5              + \dots + 0 = -2.39208984375
\end{align}

\section{Ejercicio 5}

Demostrar que si $\gamma = 0$ y $R_t = R(S_t)$ (la recompensa depende del estado), entonces
$$ \nu(s) = \sum_{s' \in S} P(S_{t+1} = s'| S_t = s) R(S_{t+1})$$.

\subsection{Resolución}

Si $\gamma = 0$ entonces el retorno es:
$$ G_t = R_{t+1} $$

Luego, el valor de un estado es:
$$ \nu(s) = E[G_t | S_t = s] = E[R_{t+1} | S_t = s] $$

Y por definición de esperanza condicional:
$$ \nu(s) = \sum_{s' \in S} P(S_{t+1} = s'| S_t = s) R(S_{t+1})$$

\section{Ejercicio 6}

Demostrar la ecuación de Bellman para MRPs con N estados
$$ \nu = R + \gamma P \nu $$

donde $v \in \mathbb{R}^N$ es el vector de valores de los estados, $r \in \mathbb{R}^N$ es el vector de recompensas medias 
partiendo de cada estado, $P \in \mathbb{R}^{N \times N}$ es la matriz de probabilidades de transiciones y 
$\gamma$ es el factor de descuento.

\subsection{Resolución}

Recordando la definición de valor de un estado:
$$ \nu(s) = E[G_t | S_t = s] $$

Si $\gamma \neq 0$ se puede reescribir como:
$$ \nu(s) = E[R_{t+1} + \gamma G_{t+1} | S_t = s] $$

Distribuyendo la esperanza:
$$ \nu(s) = E[R_{t+1} | S_t = s] + \gamma E[G_{t+1} | S_t = s] $$

para un estado $s$ específico:
$$ E[R_{t+1} | S_t = s] = \sum_{s'} P_{ss'} r_{s'} $$
donde $r_{s'}$ es la recompensa media al transicionar al estado $s'$.

A su vez, expresamos $\gamma E[G_{t+1} | S_t = s]$ como:
$$ \gamma E[G_{t+1} | S_t = s] = \gamma \sum_{s'} P_{ss'} \nu(s') $$

Finalmente, juntando todo:
$$ \nu(s) = \sum_{s'} P_{ss'} r_{s'} + \gamma \sum_{s'} P_{ss'} \nu(s') $$

factorizamos:
$$ \nu(s) = \sum_{s'} P_{ss'} (r_{s'} + \gamma \nu(s')) $$

Si expresamos esto en forma matricial, obtenemos:
$$ \nu = R + \gamma P \nu $$

\section{Ejercicio 7} % TODO

[Programación] Resuelva la ecuación de Bellman para el MRP del ejercicio 4 por los siguientes métodos
para $\gamma = 0,0.5,0.9$:

\begin{enumerate}
    \item Iterando $\nu_{n+1} = r + \gamma P \nu_n$
    \item Usando algún método iterativo para resolver sistemas lineales, por ejemplo, usando eliminación Gaussiana (via numpy.linalg.solve())
    \item Calculando la inversa de la matriz $ (I - \gamma P)$
\end{enumerate}

\subsection{Resolución}

% python code
% python code
\begin{lstlisting}[language=Python]
import numpy as np

# matriz de transicion
MATRIZ = np.array(    
    [[0,.5,0 ,0,0,.5,0], # C1
    [0,0 ,.8,0,0,0 ,.2], # C2
    [0,0 ,0 ,.6,.4,0 ,0], # C3
    [0,0 ,0 ,0,0 ,0 ,1], # PASS
    [.2,.4,.4,0,0 ,0 ,0], # Pub
    [.1,0 ,0,0 ,0,.9,0], # FB
    [0,0 ,0,0 ,0,0 ,1]] # Sleep
)

# recompensas de cada estado
R_C1 = -2
R_C2 = -2
R_C3 = -2
R_PASS = 10
R_SLEEP = 0
R_FB = -1
R_PUB = 1
GAMMA = .5

for gamma in [0,0.5,0.9]:
    I = np.eye(7)
    R = np.array([R_C1,R_C2,R_C3,R_PASS,R_PUB,R_FB,R_SLEEP])
    
    print("gamma =", gamma)
    # Iterando
    v = np.zeros(7)
    for _ in range(100):
        v = R + gamma * MATRIZ @ v
    print("\tIterando", np.round(v,2))
    
    # Eliminacion Gaussiana
    A = I - gamma * MATRIZ
    v = np.linalg.solve(A,R)
    print("\tGauss\t", np.round(v,2))
    
    # Inversa
    A = I - gamma * MATRIZ
    A_inv = np.linalg.inv(A)
    v = A_inv @ R
    print("\tInversa\t", np.round(v,2))
\end{lstlisting}

% output
\begin{lstlisting}[language=Python]
gamma = 0
	Iterando [-2. -2. -2. 10.  1. -1.  0.]
	Gauss  	 [-2. -2. -2. 10.  1. -1.  0.]
	Inversa	 [-2. -2. -2. 10.  1. -1.  0.]
gamma = 0.5
	Iterando [-2.91 -1.55  1.12 10.    0.62 -2.08  0.  ]
	Gauss  	 [-2.91 -1.55  1.12 10.    0.62 -2.08  0.  ]
	Inversa	 [-2.91 -1.55  1.12 10.    0.62 -2.08  0.  ]
gamma = 0.9
	Iterando [-5.01  0.94  4.09 10.    1.91 -7.64  0.  ]
	Gauss  	 [-5.01  0.94  4.09 10.    1.91 -7.64  0.  ]
	Inversa	 [-5.01  0.94  4.09 10.    1.91 -7.64  0.  ]
\end{lstlisting}

\section{Ejercicio 8}

Demostrar que, dada una política estocástica $\pi(a|s)$, la función de valor del par
estado-acción puede escribirse como
$$ \nu_\pi(s) = \sum_{s',r} p(s',r|s,a)[r + \gamma \nu_\pi(s')] $$

\subsection{Resolución}

Recordando la definición de la función de valor del par estado-acción:
$$ \nu_\pi(s) = E[G_t | S_t = s] $$

Si $\gamma \neq 0$ se puede reescribir como:

$$ \nu_\pi(s) = E[R_{t+1} + \gamma G_{t+1} | S_t = s] $$

Descomponemos la esperanza condicionando en todas las posibles acciones:

$$ \nu_\pi(s) = \sum_a E_\pi[G_t|S_t = s,A_t = a] * P(A_t = a|S_t = s) $$

Sabemos que $P(A_t = a|S_t = s) = \pi(a|s)$, y que $E_\pi[G_t|S_t = s,A_t = a] = q_\pi(a|s)$, por lo que:
$$ \nu_\pi(s) = \sum_a \pi(a|s) q_\pi(a|s) $$

\section{Ejercicio 9}

Demostrar que, dada una política estocástica $\pi(a|s)$, la función de valor del par
estado-acción puede escribirse como
$$ q_\pi(s,a) = \sum_{s',r} p(s',r|s,a)[r + \gamma \nu_\pi(s')] $$

\subsection{Resolución}

Sabemos que la función de valor del par estado-acción es:
$$ q_\pi(s,a) = E[G_t | S_t = s, A_t = a] $$

Remplazamos $G_t$ por $R_{t+1} + \gamma G_{t+1}$:
$$ q_\pi(s,a) = E[R_{t+1} + \gamma G_{t+1} | S_t = s, A_t = a] $$

Distribuyendo la esperanza:
$$ q_\pi(s,a) = E[R_{t+1} | S_t = s, A_t = a] + \gamma E[G_{t+1} | S_t = s, A_t = a] $$

Descomponemos la esperanza en todas las posibles recompensas y estados para
el primer término:
$$ E[R_{t+1} | S_t = s, A_t = a] = \sum_r p(r|s,a)   \stackrel{\text{bayes}}{=} \sum_{s',r} p(s',r|s,a) r $$

En el segundo término:
$$ E[G_{t+1} | S_t = s, A_t = a] = \sum_g g p(g|s,a)$$

Descomponemos $p(g|s,a)$ en todos los posibles estados siguientes:
$$ p(g|s,a) = \sum_{s'} p(g|s',s,a) p(s'|s,a) $$
Entonces queda:
$$ \sum_g g \sum_{s'} p(g|s',s,a) p(s'|s,a) $$
Sabemos que por la propiedad de Markov $p(g|s',s,a) = p(g|s')$:
$$ \sum_g g \sum_{s'} p(g|s') p(s'|s,a) $$
Además, $\nu_\pi(s') = \sum_g g p(g|s')$, por lo que:
$$ \sum_{s'} \nu_\pi(s') p(s'|s,a) $$

Remplazamos ambos términos en la ecuación final:
$$ q_\pi(s,a) = \sum_{s',r} p(s',r|s,a) r + \gamma \sum_{s'} \nu_\pi(s') p(s'|s,a) $$
Tomamos factor común:
$$ q_\pi(s,a) = \sum_{s',r} p(s',r|s,a) [r + \gamma \nu_\pi(s')] $$

\section{Ejercicio 10}

Demostrar que la función de valor de estado óptima es:
$$ \nu_*(s) = \max_a q_*(s,a) = \max_a \sum_{s',r} p(s',r|s,a)[r + \gamma \nu_*(s')] $$

\subsection{Resolución}

Recordamos que:

$$ \nu(s) = E_\pi[G_t | S_t = s] $$
$$ q_\pi(s,a) = E_\pi[G_t | S_t = s, A_t = a] $$

y por el ejercicio 9:
$$ q_\pi(s,a) = \sum_{s',r} p(s',r|s,a)[r + \gamma \nu_\pi(s')] $$

A $q_*(s,a)$ la podemos expresar como:

$$ q_*(s,a) = \max_\pi q_\pi(s,a) = \max_\pi \sum_{s',r} p(s',r|s,a)[r + \gamma \nu_\pi(s')] $$

$$ = \sum_{s',r} p(s',r|s,a)[r + \gamma \max_\pi \nu_\pi(s')] $$
$$ = \sum_{s',r} p(s',r|s,a)[r + \gamma \nu_*(s')] $$

Luego, la función de valor de estado óptima es:
$$ \nu_*(s) = \max_a q_*(s,a) = \max_a \sum_{s',r} p(s',r|s,a)[r + \gamma \nu_*(s')] $$

\section{Ejercicio 11}

Demostrar que la función de valor estado-acción verifica la siguiente ecuación:

$$ q_\pi(s,a) = \sum_{s',r} p(s',r|s,a)[r + \gamma \sum_{a'} \pi(a'|s') q_\pi(s',a')] $$

\subsection{Resolución}

Recordamos la definición de la función de valor del par estado-acción:
$$ q_\pi(s,a) = E[G_t | S_t = s, A_t = a]  = E_\pi[R_{t+1}+\gamma \sum_{k=0}^{\infty} \gamma^k R_{t+k+2}|s,a]$$

Distribuyendo la esperanza:

$$ q_\pi(s,a) = E[R_{t+1}|s,a] + \gamma E_\pi[\sum_{k=0}^{\infty} \gamma^k R_{t+k+2}|s,a] $$

Reformulamos el primer término:
$$ E[R_{t+1}|s,a] = \sum_{s',r} p(s',r|s,a) r $$

Reformulando el segundo término, tomando $G_{t+1} = \sum_{k=0}^{\infty} \gamma^k R_{t+k+2}$:
$$ E_\pi[G_{t+1}|s,a] = \sum_{g'} g' p(g'|s,a) $$
Descomponemos $p(g'|s,a)$ en todos los posibles estados siguientes:
$$ = \sum_{g'} g' \sum_{s',a'} p(g',s',a'|s,a) $$
por teorema de Bayes:
$$ = \sum_{g'} g' \sum_{s',a'} p(g'|s',a',s,a) p(s',a'|s,a) $$
por la propiedad de Markov $p(g'|s',a',s,a) = p(g'|s',a')$:
$$ = \sum_{g'} g' \sum_{s',a'} p(g'|s',a') p(s',a'|s,a) $$
y por teorema de Bayes:
$$ = \sum_{g'} g' \sum_{s',a'} p(g'|s',a') p(s'|s,a) \pi(a'|s') $$
lo reescribimos como:
$$ = \sum_{s',a'} \pi(a'|s') p(s'|s,a) \sum_{g'} g' p(g'|s',a') $$
Sabemos que $q_\pi(s',a') = \sum_{g'} g' p(g'|s',a')$, por lo que:
$$ = \sum_{s',a'} \pi(a'|s') p(s'|s,a) q_\pi(s',a') $$
Separamos las sumatorias:
$$ = \sum_{s'} p(s'|s,a) \sum_{a'} \pi(a'|s') q_\pi(s',a') $$
introducimos la variable r:
$$ = \sum_{s',r} p(s',r|s,a) \sum_{a'} \pi(a'|s') q_\pi(s',a') $$

Finalmente, juntando todo:
$$ q_\pi(s,a) = \sum_{s',r} p(s',r|s,a) r + \gamma \sum_{s',r} p(s',r|s,a) \sum_{a'} \pi(a'|s') q_\pi(s',a') $$

Tomamos factor común:
$$ q_\pi(s,a) = \sum_{s',r} p(s',r|s,a) [r + \gamma \sum_{a'} \pi(a'|s') q_\pi(s',a')] $$

\section{Ejercicio 12}

Demostrar que la función de valor de estado-acción óptima es:

$$ q_*(s,a) = \sum_{s',r} p(s',r|s,a) [r + \gamma \max_{a'} q_*(s',a')] $$

\subsection{Resolución}

Definimos $q_*(s,a)$ como:
$$ q_*(s,a) = \max_\pi q_\pi(s,a)$$

Entonces queremos demostrar que:
$$ \max_\pi [\sum_{a'}q_\pi(s',a') \pi(a'|s')] $$

El máximo se obtiene con la política greedy, es decir, eligiendo la acción que maximiza $q_\pi(s',a')$:
$$ \max_\pi [\sum_{a'}q_\pi(s',a') \pi(a'|s')] = \max_{a'} q_\pi(s',a') $$

Luego, por el ejercicio 11:
$$ q_*(s,a) = \sum_{s',r} p(s',r|s,a) [r + \gamma \max_{a'} q_*(s',a')] $$

\end{document}
